\section{Related Works}
\label{sec:related_works}

\subsection{Demographic Face Recognition}
\label{sec:demographic_face_recognition}

The development of facial demographic attribute analysis originated from the utilization of local geometric descriptors and handcrafted features, such as Local Binary Patterns and Gabor filters, progressing toward deep CNN-based representations. Local region-based approaches leverage the geometry of the middle-face region to recognize ethnicity while suppressing background image interference~\cite{ref11}. However, classical convolutional architectures still face challenges in robustness against extreme illumination fluctuations and cross-racial performance disparities known as the Other-Race Effect. Increasing network depth has been demonstrated to broaden the discriminative capacity of visual features in distinguishing heterogeneous populations~\cite{ref13}. In addition to deepening layers, parallel feature-sharing mechanisms with visual attention modules have been developed to capture correlations between facial attributes~\cite{ref16}, reinforcing the transition toward representations that are more adaptive to facial morphological diversity.

Following feature extraction, the determination of downstream classifiers plays a crucial role in mapping latent representations into demographic decision spaces. Conventional models generally rely on end-to-end SoftMax layers to generate class probabilities~\cite{ref11}, \cite{ref18}. However, non-linear mappings such as Kernel Extreme Learning Machine (KELM) and SVM on frozen representations offer robust separation boundaries without backbone weight updates~\cite{ref13}, \cite{ref12}. The mathematical formulation of the KELM decision function is modeled in \eqref{eq:1}:
\begin{equation}
f(\mathbf{x}) = \mathbf{h}(\mathbf{x})\boldsymbol{\beta} = \sum_{i=1}^M \beta_i K(\mathbf{x}, \mathbf{x}_i)
\label{eq:1}
\end{equation}
where $\mathbf{h}(\mathbf{x})$ denotes the latent feature mapping vector, $\boldsymbol{\beta} = [\beta_1, \dots, \beta_M]^T$ represents the optimized weights, $K(\mathbf{x}, \mathbf{x}_i)$ is the non-linear kernel function with respect to support samples $\mathbf{x}_i$, and $M$ denotes the number of representation bases. This kernel mapping enables effective non-linear decision boundary separation in high-dimensional spaces without iterative gradient descent optimization.

\subsection{Vision Transformers in Face Analysis}
\label{sec:vision_transformers_face}

The application of ViT in facial biometrics introduces global modeling that overcomes the local inductive bias of CNNs~\cite{ref14}. The input image is partitioned into $N$ patches of resolution $P \times P$ with $C$ color channels, and subsequently projected into the latent dimension via formulation \eqref{eq:2}:
\begin{equation}
\mathbf{x}_p \in \mathbb{R}^{N \times (P^2 \cdot C)}, \quad \mathbf{z}_0 = [\mathbf{x}_{\text{class}}; \, \mathbf{x}_p^1\mathbf{E}; \dots; \, \mathbf{x}_p^N\mathbf{E}] + \mathbf{E}_{\text{pos}}
\label{eq:2}
\end{equation}
where $N = HW/P^2$ denotes the number of patches, $P$ is the patch size, $C$ represents the color channels, $\mathbf{E} \in \mathbb{R}^{(P^2 \cdot C) \times D}$ is the projection matrix, $\mathbf{x}_{\text{class}}$ is the class token, and $\mathbf{E}_{\text{pos}}$ denotes the positional encoding. Global inter-patch spatial interactions are modeled through the MHSA mechanism according to \eqref{eq:3}:
\begin{equation}
\text{Attention}(\mathbf{Q}, \mathbf{K}, \mathbf{V}) = \text{softmax}\left(\frac{\mathbf{Q}\mathbf{K}^T}{\sqrt{d_k}}\right)\mathbf{V}
\label{eq:3}
\end{equation}
with $\mathbf{Q}$, $\mathbf{K}$, and $\mathbf{V}$ representing the Query, Key, and Value matrices, and $\sqrt{d_k}$ denoting the scaling factor. This attention mechanism has been effectively applied to facial gender classification~\cite{ref17}, while MaxViT enriches representational hierarchy through block and grid attention~\cite{ref18}.

The global representation characteristics produced by transformer architectures have prompted thorough investigations into the propagation of demographic bias in computer vision systems. Evaluation on balanced data corpora provides controlled experimental conditions to compare performance disparities across subgroups without the distortion of class sample size imbalances. Although data distributions are equalized, such balance does not automatically eliminate the social biases inherent in the model. Investigations into latent representations demonstrate that transformer architectures can reproduce and amplify social biases originating from massive pre-trained corpora~\cite{ref14}. As a mitigation measure for such disparities, a generative perspective approach based on StyleGAN2 has been utilized to enrich phenotypic diversity across demographic groups~\cite{ref15}. The integration of these methods demonstrates that auditing representation fairness and aligning data distributions are crucial in the design of modern facial biometric models.

\subsection{Multi-Domain Feature Fusion}
\label{sec:multidomain_feature_fusion}

The limitations of single-domain representations have driven the development of multi-task learning and cross-representation fusion. Multi-task learning has proven effective in jointly mapping the correlations between age and gender estimation~\cite{ref12}, whereas representation-sharing schemes facilitate the simultaneous extraction of various affective attributes~\cite{ref16}. Furthermore, the integration of ViT-based latent features demonstrates the advantages of combining biometric domains with emotional expression~\cite{ref19}, as well as biological age~\cite{ref20}. Such orthogonal multi-domain representation fusion is formulated through the direct sum operation in \eqref{eq:4}:
\begin{equation}
\mathbf{z}_{\text{fused}} = \mathbf{z}_{\text{face}} \oplus \mathbf{z}_{\text{emotion}} \oplus \mathbf{z}_{\text{age}} \in \mathbb{R}^{D_{\text{total}}}
\label{eq:4}
\end{equation}
where the operator $\oplus$ denotes the direct concatenation of equally weighted feature vectors from the biometric domain $\mathbf{z}_{\text{face}}$, affective expression $\mathbf{z}_{\text{emotion}}$, and biological age $\mathbf{z}_{\text{age}}$, with a combined dimension of $D_{\text{total}} = 2{,}304$. This integration comprehensively fuses static identity cues, dynamic expression variations, and natural aging patterns.

Considering the methodological limitations of prior studies, this study establishes its novelty through an integrated tri-domain fusion combining facial biometric representations, affective emotional expressions, and biological age from pre-trained ViT models with frozen backbone weights. Distinct from prior literature that evaluated performance on image-level random partitioning schemes susceptible to identity leakage, this framework is rigorously evaluated under a Strict Subject-Disjoint Partitioning protocol to ensure generalization validity on unseen subjects. Furthermore, the contributions of this study encompass a systematic exploration of seven feature representation ablation schemes across four classical machine learning classifier families, leakage-free fold-isolated preprocessing optimization, and a formal inferential characterization of the non-linear representation saturation phenomenon versus linear complementarity.
